% Mathematical source: IUT_IDENTITY_LOG_LINK_LOCAL_MEMBERSHIP_2026_08_31.md
% Accepted report SHA256:
% 6ca3f92988be870df06b3536d9f1f6b598e9ed1f87cbd2feef5d8785e3f6d3d9
% This input does not add a Lean formalization of source reconstruction.
\section{A synchronized base branch of the theta-pilot images}
\label{sec:uniform-identity-membership}

We now place the local points of
Section~\ref{sec:uniform-theta-points} in a specified base branch
of the possible-image construction in
\cite[Corollary~3.12]{IUT3}. The construction uses one concrete
global theta-pilot, two tagged copies of its Hodge theater,
and the log-link representative induced by their single global
copy-identity. The resulting local inclusion concerns the raw
union of possible images in one vertically coric column.
It does not identify the complete global family with the
one-prime arithmetic bundles of
Section~\ref{sec:uniform-arithmetic-bundles}.

We use the reconstruction results of the cited sources, with
their original objects and variance conventions. We do not
reprove all those results, or assume the global multiradial
and horizontal compatibility assertions of
\cite[Theorem~3.11]{IUT3} in order to establish the inclusion
below. References to \cite{IUT1,IUT3}, \cite{IUT2}, and
\cite{MochizukiAbsTopIII2015} use respectively the archived
May 2020, December 2020, and November 2015 author versions.

\subsection{Fixed initial data and the standard log-field}

Fix a member of Corollary~\ref{cor:itd-powerfree-family}.
Its global initial theta data, including its core, auxiliary
cover, distinguished quotient element and place section, have
already been constructed. At its distinguished rational
prime \(p\), retain \eqref{eq:utp-field}: in particular,
\[
 e=15\ell,\qquad d=2e,\qquad h=(\ell-1)/2,\qquad
 I=p^{-1}\log\OO_E^\times=\beta^{1-e}\OO_E,\qquad
 r_0^{2\ell}=q=b_0^2.
\]
Here \(v_p(b_0)=2\), \(E/\Q_p\) is Galois with residue
degree two, \(e\le p-2\), and
\(\mu_{2\ell}\subset K_0=\Q_p(\mu_{30\ell})\).
For \(1\le j\le h\), put
\begin{equation}\label{eq:uim-point-data}
 \begin{gathered}
 m_j=j+1,\qquad a_j=r_0^{j^2},\qquad
 x_j=\zeta_j a_j,\quad \zeta_j\in\mu_{2\ell},\\
 \mathcal T_{m_j}=E^{\otimes_{\Q_p}m_j},\qquad
 Z_{j,\zeta_j}=1^{\otimes j}\otimes x_j .
 \end{gathered}
\end{equation}
The distinguished last slot is the label-\(j\) slot in the
source procession; it is not an extra factor.

The field of moduli is \(F_{\rm mod}=\Q\).
The source's selected place set \(V\) maps bijectively to
the places of \(F_{\rm mod}\), so it has exactly one
selected \(v\) above \(p\). This is not a claim that the
global torsion field has only one prime above \(p\).
It means that the finite local packet at this rational
prime has no other selected direct summands:
its marked carrier is \(\mathcal T_{m_j}\).

Let \(E^-\) denote the original predecessor field.
Definition~1.1(i) of \cite[pp.~23--25]{IUT3} equips the
perfection of its multiplicative unit group with a
\emph{new field structure}, denoted here by \(L^-\).
In the concrete model, standard logarithm gives a
unital field isomorphism
\begin{equation}\label{eq:uim-log-field}
                 \mathcal L:L^-\xrightarrow{\sim}E.
\end{equation}
It transports both field operations from \(E\);
the new multiplication is not the predecessor's
old multiplication of units.

For \(x\in E\), choose \(N\) sufficiently large and write
\begin{equation}\label{eq:uim-perfect-unit}
       y(x)=p^{-N}[\exp(p^Nx)]\in L^- .
\end{equation}
The bracket is written additively in the old unit
perfection. The equality
\(\exp(p^{N+1}x)=\exp(p^Nx)^p\) proves independence of
sufficiently large \(N\), and
\(\mathcal L(y(x))=x\).
For \(x\in I\), one may take \(N=1\), since
\(pI=\mathfrak m_E\) and \(1/e>1/(p-1)\).

In the chart \eqref{eq:uim-log-field}, the pre-log-shell
has coordinates \(\log\OO_E^\times\); the canonical
shell is its multiple \(p^{-1}\log\OO_E^\times=I\).
This division by \(p\) does not change the field
multiplication or its unit. The new field unit is
\(y(1)\), of coordinate \(1\), whereas the old
multiplicative unit \(1\in(E^-)^\times\) has coordinate
zero after logarithm. Thus the background entries of
\eqref{eq:uim-point-data} are the new field units.

For comparison, changing the \emph{entire field chart}
to \(u=x/p\) would instead give
\begin{equation}\label{eq:uim-rescaled-field}
 \begin{gathered}
 u\star v=p\,uv,\qquad 1_\rho=1/p,\qquad
 \Q_p\ni t\longmapsto t/p,\\
 v_\rho(u)=v_E(pu),\qquad
 \OO_\rho=\OO_E/p,\qquad I_\rho=I/p .
 \end{gathered}
\end{equation}
The equation \(a^n=b\) would become
\(p^{n-1}a_\rho^n=b_\rho\). In an \(m\)-fold packet,
both the tensor coordinates and the reference integer
module would be transported by \(p^{-m}\).
This is different from changing only the outputs of a
fixed unit-cohomology source from
\(\rho=p^{-1}\log_{\rm BK}^{\rm std}\) to
\(\log_{\rm BK}^{\rm std}\), which gives the factor
\(p^m\) in \eqref{eq:utp-standard-scale}.
Neither rescaling is made in \eqref{eq:uim-log-field}.

\subsection{A global identity and three distinct cyclotomes}

Take distinct tagged copies
\(\mathcal H_{n,r}\), \(n,r\in\Z\), of the fixed
concrete Hodge theater. Retain the full vertical and
horizontal poly-arrows required by
\cite[Definition~3.8(iii), pp.~113--114]{IUT3}.
Inside the vertical arrow
\(\mathcal H_{0,-1}\to\mathcal H_{0,0}\), select the
representative induced by the global copy-identity
\begin{equation}\label{eq:uim-global-identity}
 \Xi:\mathcal H^D_{0,-1}\xrightarrow{\sim}
                         \mathcal H^D_{0,0}.
\end{equation}
Selecting a representative does not replace a full
poly-arrow by a singleton.

\begin{proposition}[Synchronized standard-log representatives]
\label{prop:uim-synchronized}
The single isomorphism \eqref{eq:uim-global-identity}
gives, at every finite place and every labeled
prime-strip occurrence, the log-field map
\begin{equation}\label{eq:uim-local-identity}
       \lambda_v=\operatorname{copy}_v\circ\mathcal L_v:
           L_{0,-1,v}\xrightarrow{\sim}E_{0,0,v}.
\end{equation}
These are simultaneously allowed representatives
of the theater log-link. At the distinguished place,
the pullbacks of \(x_j\) and of each background field
unit have standard-log coordinates \(x_j\) and \(1\).
\end{proposition}
\begin{proof}
Definition~1.1(i) of \cite{IUT3} constructs the natural
isomorphism from the tautological log-Frobenioid to
the original Frobenioid. The group acts compatibly
on all the monoids in that construction.
Remark~1.1.2(i), pp.~29--30, identifies its underlying
local fundamental groups by the tautological identity.
Composing this natural isomorphism with the copy map
therefore covers the D-prime-strip isomorphism
induced by \(\Xi\).

The map from F-prime-strip isomorphisms to their
D-prime-strip isomorphisms is bijective by
\cite[Corollary~5.3(ii), p.~144]{IUT1}.
Consequently the displayed local map is the unique
lift used in
\cite[Proposition~1.3(i), p.~42]{IUT3}.
That proposition applies these lifts to the one
\(\Xi\) across every prime-strip of both bridges;
Remark~1.3.1, p.~43, makes the synchronization
requirement explicit. The construction at the
archimedean constituents uses the corresponding
standard tautological isomorphism.

The concrete marked theater has the synchronized
constant-monoid identifications of
\cite[Corollary~4.6(iii), p.~138]{IUT2}.
Its copy-identity neither permutes the labels nor
changes these identifications. In those fixed
coordinates \eqref{eq:uim-local-identity} is the
identity field map. The asserted pullbacks follow.
\end{proof}

The proposition is about this \(\Xi\), not every
theater isomorphism. An arbitrary such isomorphism
may change label markings. A unital local field
isomorphism fixing \(q\in\Q_p\) preserves the coset
defined by \(X^{2\ell}=q^{j^2}\), but that observation
alone does not provide a synchronized global
representative.

It is essential to distinguish the three coefficient
objects
\[
 \mu_{\rm prev},\qquad \mu_{\log},\qquad
 \mu_{\rm current},
\]
arising respectively from multiplication in the
old field \(E^-\), the new log-field \(L^-\), and
the current field \(E^+\).
The old \emph{linear unit-Kummer coordinate} is
\begin{equation}\label{eq:uim-linear-kummer}
 \kappa^-_{\rm lin}(x)
   =p^{-N}\operatorname{Kum}_{\mu_{\rm prev}}
                   (\exp(p^Nx))
       \in H^1_f(G_E,\Q_p(1)_{\rm prev}).
\end{equation}
It is the inverse standard Bloch--Kato logarithm
on the unit subspace. On \(I\) it takes values in
\(p^{-1}H^1_f(G_E,\Z_p(1)_{\rm prev})\).

There is separately the new multiplicative Kummer map
\[
 (L^-)^{\times,\log}\longrightarrow
       H^1(G_{L^-},\Q_p(1)_{\log}).
\]
The field isomorphism \(\lambda\) of
\eqref{eq:uim-local-identity} gives the usual
field-theoretic Kummer comparison
\begin{equation}\label{eq:uim-new-kummer}
 \operatorname{Kum}_{\mu_{\log}}(y(x))
       \longmapsto
 \operatorname{Kum}_{\mu_{\rm current}}(\lambda(y(x))).
\end{equation}
Its group and coefficient maps are those induced by
that field isomorphism. The operation
\(y\mapsto\operatorname{Kum}_{\mu_{\log}}(y)\)
is not linear in the old unit-H1 coordinate
\eqref{eq:uim-linear-kummer}.
The new multiplication used here is already part
of \cite[Definition~1.1(i), p.~24]{IUT3}, and its
monoid, units and action are used explicitly in
\cite[Proposition~3.4(i), p.~102]{IUT3}.

Indeed, a same-group integral-coefficient H1 arrow
cannot identify the two Kummer maps just displayed.
If \(v_E(x)>0\), the raw class
\(\operatorname{Kum}_{\mu_{\rm current}}(x)\) has
nonzero valuation component, while
\(\kappa^-_{\rm lin}(x)\) has valuation component
zero. The unit subspace is the annihilator of
unramified characters under local Tate duality,
so a Galois-induced isomorphism with compatible
integral Tate coefficients preserves it.
The comparison \eqref{eq:uim-new-kummer} concerns
the new multiplicative structure and does not
assert this impossible old-H1 identification.

No global exponential is being used. The element
\(\exp(p^Nx)\) in \eqref{eq:uim-perfect-unit}
need only belong to the completion \(E\).
It is not a global rational-function torsor
generator. The source global localization is
pulled back along the globally synchronized
local isomorphisms
\cite[Proposition~3.3]{IUT3}.
The same distinction is explicit in
\cite[Definition~5.4(ii), pp.~125--126]{MochizukiAbsTopIII2015}:
the global log construction re-embeds the
global object through its natural local
isomorphisms and keeps the underlying global
Galois theater. Its Remark~5.4.1, p.~129,
warns that the local logarithm diagrams
do not extend as a logarithm on the global
number field.

\subsection{The canonical core coordinate and the preceding layer}

Let \(\mathcal C_E\) denote the finite
Galois-invariant part of the canonical module
\(k^\sim(G_E)\) in
\cite[Proposition~5.8(ii), p.~139]{MochizukiAbsTopIII2015}.
Use the local reciprocity convention of that
unit reconstruction,
\(\operatorname{rec}_E:E^\times\to G_E^{\rm ab}\).
The concrete model of \(\mathcal C_E\) is the
rationalized old unit group with its finite
torsion removed.

\begin{proposition}[The specified unit comparison]
\label{prop:uim-core-coordinate}
There is a standard comparison
\begin{equation}\label{eq:uim-sigma}
 \sigma_E:E\xrightarrow{\sim}\mathcal C_E,\qquad
 \sigma_E(x)
     =p^{-N}[\operatorname{rec}_E(\exp(p^Nx))],
\end{equation}
and the source shell is \(\sigma_E(I)\).
For the unit Kummer comparison
\begin{equation}\label{eq:uim-eta}
 \eta_E:H^1_f(G_E,\Q_p(1)_{\rm prev})
          \xrightarrow{\sim}\mathcal C_E,\qquad
 \eta_E(\operatorname{Kum}(u))
          =[\operatorname{rec}_E(u)],
\end{equation}
one has
\begin{equation}\label{eq:uim-sigma-kappa}
                    \sigma_E=\eta_E\kappa^-_{\rm lin}.
\end{equation}
\end{proposition}
\begin{proof}
The brackets in \eqref{eq:uim-sigma} refer to
the unit-perfection module. Independence of
sufficiently large \(N\) follows from the
exponential identity used for
\eqref{eq:uim-perfect-unit}. Standard logarithm
on a sufficiently small unit subgroup proves
additivity, \(\Q_p\)-linearity and bijectivity.
The actual unit image is the pre-log-shell,
and its multiple \(p^{-1}\) is the shell.
This gives \(\sigma_E(I)\), with no further
factor \(p^{-1}\) in \eqref{eq:uim-sigma}.

The Kummer sequence and Hilbert 90 identify
the rationalized torsion-free unit group
with \(H^1_f\), so \eqref{eq:uim-eta} is
well defined. Evaluating it on
\(p^{-N}\operatorname{Kum}(\exp(p^Nx))\)
proves \eqref{eq:uim-sigma-kappa}.

This is the particular comparison specified
by the reconstruction source.
Parts~(c) and~(d) of
\cite[Corollary~5.10(iv), p.~148]{MochizukiAbsTopIII2015}
give the natural
comparison from the field-theoretic shell
to the canonical mono-analytic shell,
compatible with the shells and volumes.
Its proof on p.~149 compares the base field
inside abelianizations with its Kummer
reconstruction by the fundamental-class
and cyclotomic isomorphisms of
Corollary~1.10(a),(c). The proof identifies
the result as the reciprocity map of
local class field theory. On a concrete
unit \(u\), this is
\eqref{eq:uim-eta}. Thus the comparison
used here is not an arbitrary linear
isomorphism chosen after the point.
\end{proof}

Write \(k_-^\times,k_+^\times\) for the
original constant-monoid reconstruction
maps of
\cite[Corollary~4.6(iii)]{IUT2}, from the
two copied theaters to the common core.
In the fixed concrete model they reconstruct
the same labeled field. Before forgetting
multiplication, the relevant diagram is
\begin{equation}\label{eq:uim-field-square}
 \begin{array}{ccc}
 L^-&\xrightarrow{\lambda}&E^+\\
 {\scriptstyle\log(k_-^\times)}\downarrow
       &&\downarrow{\scriptstyle k_+^\times}\\
 L^{\rm c}&\xrightarrow{\lambda_{\rm c}}&E^{\rm c}.
 \end{array}
\end{equation}
The lower map is the corresponding tautological
standard-log representative.
Both paths send \(y(x)\) to the constant-field
element of coordinate \(x\). To check this directly,
first take a sufficiently small old unit:
its constant-monoid reconstruction is its
actual reconstructed constant, and standard
logarithm commutes term by term with this
copy-identity field reconstruction.
Division by \(p^N\) extends the equality to
every \(y(x)\). The left arrow is formed from
the old-unit map and then the new log-field
structure; its value in the core additive
carrier is \(\sigma_E(x)\).
The right arrow is a current multiplicative
constant-monoid map. These are not two
versions of the same H1 arrow.

Let
\[
 \mathcal C_{j,p}
   =I^\Q(S^\pm_{j+1};{}^{0,\circ}D^\perp_p),
 \qquad
 \iota_j=\sigma_E^{\otimes m_j}:
           \mathcal T_{m_j}\xrightarrow{\sim}\mathcal C_{j,p}.
\]
The superscript \(\perp\) here denotes the
source's mono-analytic prime-strip, not an
orthogonal complement. Its identification with
the tensor of \(\mathcal C_E\) is the natural
one in \cite[Proposition~3.2(i),(ii), pp.~97--99]{IUT3}.
Denote the label-\(j\) components of the
split LGP monoids at level \(0\) and at the
core by \(\mathcal S_j^+\) and
\(\mathcal S_j^{\rm c}\). Applying the
distinguished-slot recipe to
\eqref{eq:uim-field-square} gives
\begin{equation}\label{eq:uim-packet-square}
 \begin{array}{ccc}
 \mathcal S_j^+&\lhook\joinrel\longrightarrow&
                                    \mathcal T_{m_j}\\
 {\scriptstyle k_{{\rm LGP},0}}\downarrow
       &&\downarrow{\scriptstyle\iota_j}\\
 \mathcal S_j^{\rm c}&\lhook\joinrel\longrightarrow&
                                    \mathcal C_{j,p}.
 \end{array}
\end{equation}
The upper right is written in the standard-log
coordinates of the \emph{preceding} field \(L^-\).
Before choosing coordinates the right-hand map
is \(\iota_j\mathcal L^{\otimes m_j}\).

Here the upper inclusion is
\cite[Proposition~3.4(ii), pp.~102--103]{IUT3}:
pull the current Gaussian monoid back to the
preceding log-field, then tensor with the
new units in the other slots.
The lower inclusion is the same recipe on
the reconstructed constant monoids.
Proposition~3.5(i), pp.~103--104, explicitly
uses the comparisons at \emph{both}
indices \(r'=r,r-1\).
The commutativity of
\eqref{eq:uim-field-square}, followed by
this tensor operation, verifies
\eqref{eq:uim-packet-square} in the selected
concrete base representative.
It is not obtained by replacing the
preceding carrier with the current H1 space.

The same diagram at all places and the
copy-identity of the global field give
the corresponding base fractional-ideal
comparison. Both recipes use the same
local ideals and the same global
rational-function monoid acting on them.
This is the construction recorded in
\cite[Proposition~3.10(i), pp.~147--148]{IUT3},
again with both adjacent indices.
We use this fixed representative, not the
later assertion of compatibility for every
vertical iterate.

\begin{proposition}[The variance of the actual core action]
\label{prop:uim-transfer}
For a continuous automorphism \(\alpha\) of \(G_E\),
let \(M_\alpha\) be its action on the unit
abelianization, expressed in logarithmic coordinates:
\begin{equation}\label{eq:uim-canonical-matrix}
 M_\alpha(\log u)
   =\log\!\left(\operatorname{rec}_E^{-1}
       \alpha^{\rm ab}(\operatorname{rec}_E(u))\right).
\end{equation}
Extend this map \(\Q_p\)-linearly to \(E\).
The actual canonical core map \(\Phi_\alpha\)
satisfies
\begin{equation}\label{eq:uim-transfer}
       \Phi_\alpha\sigma_E=\sigma_E M_\alpha^{-1},
       \qquad
       \Phi_\alpha(\sigma_E(I))=\sigma_E(I).
\end{equation}
\end{proposition}
\begin{proof}
The unit image in the abelianization is preserved
by the reconstruction algorithm.
Proposition~5.8(i) of
\cite[p.~139]{MochizukiAbsTopIII2015}
specifies its \emph{contravariant} functoriality
by transfer. For a group isomorphism this
transfer is the inverse of the induced
abelianization map; there is no index factor.
Thus
\[
 \Phi_\alpha(\sigma_E(x))
  =p^{-N}[\alpha^{-1,{\rm ab}}
       (\operatorname{rec}_E(\exp(p^Nx)))]
  =\sigma_E(M_\alpha^{-1}x).
\]
The unit image and its perfection are preserved.
Dividing the unit image by \(p\) proves the
shell assertion.
\end{proof}

For comparison, a Kummer action with a
chosen compatible integral Tate-module
isomorphism \(\delta\) is
\[
 F_{\alpha,\delta}=c_{\alpha,\delta}M_\alpha^{-1},
                 \qquad c_{\alpha,\delta}\in\Z_p^\times,
\]
by Proposition~\ref{prop:gl-variance}.
We do not silently identify this with
\eqref{eq:uim-transfer}. All pointwise
statements below use the actual canonical
core operator \(F=M_\alpha^{-1}\).
In particular, if a full Galois lift
\(\alpha_0\) has canonical matrix \(M\),
then \(\alpha_0^{-1}\) gives core matrix
\(M_{\alpha_0^{-1}}^{-1}=M\).
No additional scalar choice is required.

\subsection{Local ideal elements of one global pilot}

\begin{lemma}[Points belonging to one pilot object]
\label{lem:uim-pilot-points}
In the synchronized representative above,
there is one global theta-pilot
\(\mathcal P\), in the local fractional-ideal
realization of
\cite[Definition~3.8(i)]{IUT3}, whose
label-\(j\) ideal at the distinguished
place contains \(Z_{j,\zeta_j}\) for
every \(\zeta_j\in\mu_{2\ell}\).
\end{lemma}
\begin{proof}
Choose one standard normalized theta-evaluation
representative. Its label-\(j\) value is a
\(\mu_{2\ell}\)-multiple of \(a_j\), by
Lemma~\ref{lem:utp-evaluation} and
\cite[Remark~2.5.1(i), p.~72]{IUT2}.
The Gaussian splitting monoid is generated,
up to torsion, by this value-profile
\cite[Corollary~3.6(ii),(iii), pp.~99--100]{IUT2};
Corollary~4.6(iv), pp.~138--139, supplies
it in the chosen theater.

Proposition~\ref{prop:uim-synchronized} and
\cite[Proposition~3.4(ii)]{IUT3} pull that
generator to the preceding log-field with
the same standard coordinate.
The single-slot insertion of
\cite[Proposition~3.1(ii), p.~93]{IUT3}
tensors it with the new field units.
The local fractional ideal generated by
this element contains its generator
times the unit. Multiplication by a
root of unity does not change that
ideal, since these are units of the
new log-field. This proves the local
claim for every \(\zeta_j\).

At the other bad places choose the
prescribed splitting generators of the
same theater, and use its standard
good-place constituents. The finite
collection of bad-place local ideals
defines a global object by
\cite[Proposition~3.7(v), pp.~111--112]{IUT3}.
Definition~3.8(i), p.~112, is precisely
the theta-pilot designation for that
object. Thus all the local ideal
elements in the statement belong
to one \(\mathcal P\).
The construction does not require
the local generators to be a
single element of the global
number field.
\end{proof}

The lemma uses elements of the actual local
ideals. It need not assert that every
independent tuple \((\zeta_j)\) arises
from one strictly chosen global
theta-evaluation representative.
The ideals, and hence the global
pilot object, are unchanged by
those torsion choices.

We also record exactly which local
Galois operations are permitted.
The Ind1 object in
\cite[Theorem~3.11(i), pp.~153--154]{IUT3}
is the output mono-analytic procession
of \cite[Proposition~6.9(ii)]{IUT1}.
At the selected nonarchimedean place,
its constituent category is the
category of finite transitive
continuous \(G_E\)-sets.
Pullback along any continuous
\(\alpha\in\operatorname{Aut}(G_E)\)
is a self-equivalence, with inverse
the pullback along \(\alpha^{-1}\).
Use this equivalence at every
repeated occurrence above \(p\)
and the identity elsewhere,
keeping the place and capsule
indices fixed.

Definition~4.1(iii),(iv), p.~96,
of \cite{IUT1} makes this a
constituent prime-strip isomorphism.
Its Section~0's capsule-full
poly-isomorphisms contain the
chosen collection with the identity
index map. Definition~4.10,
pp.~119--120, makes these precisely
the representative data permitted
at the stages of the procession.
The same choice at repeated
occurrences is allowed, and
conjugation leaves the full
connecting collections unchanged.
Thus the collection occurs inside
an allowed procession
poly-automorphism, as in the
representative argument of
Section~\ref{sec:admissible-arrows}.
No extension of \(\alpha\) to the
earlier curve covering category
or global Hodge theater is imposed
by this output object.

\begin{theorem}[Membership in a specified base branch]
\label{thm:uim-base-membership}
Let \(\mathcal P\) be the fixed global
pilot of Lemma~\ref{lem:uim-pilot-points},
and let \({}^{0,\circ}U^{\rm raw}_{j,p}\)
denote the raw possible-image union
defined in \cite[Corollary~3.12,
p.~174]{IUT3}, before its holomorphic
hull is taken. For every continuous
\(\alpha\in\operatorname{Aut}(G_E)\),
put \(F=M_\alpha^{-1}\).
Then, simultaneously for all labels,
\begin{equation}\label{eq:uim-membership}
 \iota_j\!\left(F^{\otimes m_j}Z_{j,\zeta_j}\right)
       \in{}^{0,\circ}U^{\rm raw}_{j,p}
                     \qquad(1\le j\le h).
\end{equation}
These images use the fixed adjacent-layer
branch \eqref{eq:uim-packet-square},
the same local Ind1 representative
at every occurrence above \(p\),
the identity elsewhere, and identity
Ind2. The assertion is membership
for this branch, not compatibility
of all branches.
\end{theorem}
\begin{proof}
Lemma~\ref{lem:uim-pilot-points} puts
\(Z_{j,\zeta_j}\) in the local ideal
of \(\mathcal P\).
The concrete base construction
\eqref{eq:uim-packet-square},
including its fractional-ideal
realization, therefore places
\(\iota_j Z_{j,\zeta_j}\) in the
core realization of that same
pilot. The permitted representative
described above acts on this
local packet by
\(\Phi_\alpha^{\otimes m_j}\).
Proposition~\ref{prop:uim-transfer}
gives
\[
 \Phi_\alpha^{\otimes m_j}
       (\iota_jZ_{j,\zeta_j})
       =\iota_j F^{\otimes m_j}Z_{j,\zeta_j}.
\]
This is thus an element of the
image of an actual local ideal
of the pilot under an allowed
representative.

The raw set on p.~174 is the union
of such possible images.
The basic Kummer image is the
branch specified in
\cite[Proposition~3.5(ii)(a)(1),
pp.~104--105]{IUT3}; here it has
not been precomposed with a
positive log-iterate. Including
the further vertical branches
in that union does not remove
the chosen image.
This proves \eqref{eq:uim-membership}.
In particular, we have not
treated Ind3 as a group with
an identity element or used
a claimed commutativity of
every vertical Kummer map.
\end{proof}

\begin{remark}[Transport of an ideal and its test module]
\label{rem:uim-operator}
The source gives the splitting monoid
both as a subset of the packet
and as a multiplicative action
\cite[Theorem~3.11(i)(b),
pp.~153--154]{IUT3}.
If its local ideal is presented
as \(J=m_z(R)\), with \(1\in R\),
then any additive carrier
isomorphism \(\Phi\) satisfies
\begin{equation}\label{eq:uim-operator}
 \begin{gathered}
 \Phi J=(\Phi m_z\Phi^{-1})(\Phi R),\\
       (\Phi m_z\Phi^{-1})(\Phi1)=\Phi z .
 \end{gathered}
\end{equation}
Both identities follow by applying
the maps to an element of \(R\).
They require no multiplicativity
of \(\Phi\).
Using the conjugated operator
on a fixed native \(1\) or a
fixed native integer module
would be a different operation.
The proof of
Theorem~\ref{thm:uim-base-membership}
uses \(\Phi J\), with the test
module and vector transported
as in \eqref{eq:uim-operator}.
Its embedded global field is
also transported; it is not
required to be stabilized in
its old native coordinates.
\end{remark}

For general initial data with more
than one selected place above \(p\),
the ambient carrier is
\(\bigotimes_{\rm slots}
(\bigoplus_{v\mid p}E_v)\).
The point \(1^{\otimes j}\otimes x_j\)
in a specified all-\(v\) block
is then the projection of the
source's actual single-slot
point. The corresponding
statement is membership in
the projection of the full
possible-image set, with the
same global pilot image as
preimage. Zero extension of
that block does not follow.
The condition \(F_{\rm mod}=\Q\)
above is what permits the
literal membership
\eqref{eq:uim-membership}
in the whole local packet.

\subsection{An attained local lower inclusion and its boundary}

Let
\[
 \Gamma_E^{\rm can}
   =\{M_\alpha^{-1}:
          \alpha\in\operatorname{Aut}_{\rm cont}(G_E)\}.
\]
These are the canonical core
operators, expressed in the
fixed standard-log chart.
Retain \(k_j=\lfloor2j^2/\ell\rfloor+1\),
\(\kappa=1-1/e\), and the
normalized tensor integer ring
\(B_{m_j}\) of
Section~\ref{sec:uniform-theta-points}.

\begin{corollary}[An attained lower ideal in the possible images]
\label{cor:uim-canonical-hull}
For every tuple
\((\zeta_1,\ldots,\zeta_h)\in\mu_{2\ell}^h\),
there is one \(F\in\Gamma_E^{\rm can}\) such that
\begin{equation}\label{eq:uim-common-minimum}
 \begin{gathered}
 v_E(F(1))=-\kappa,\\
 v_E(F(x_j))=k_j-\kappa\qquad(1\le j\le h).
 \end{gathered}
\end{equation}
For every label and fixed torsion representative,
\begin{equation}\label{eq:uim-canonical-hull}
 \begin{split}
 \overline{\Span_{B_{m_j}}
    \{F^{\otimes m_j}Z_{j,\zeta_j}:
                                F\in\Gamma_E^{\rm can}\}}
      &=P_j,\\
 P_j&=\beta^{\,e k_j-(e-1)m_j}B_{m_j}.
 \end{split}
\end{equation}
Consequently, in this fixed local chart,
\begin{equation}\label{eq:uim-lower-inclusion}
 P_j\subseteq
 \overline{\Span_{B_{m_j}}
       \bigl(\iota_j^{-1}({}^{0,\circ}U^{\rm raw}_{j,p})\bigr)}.
\end{equation}
\end{corollary}
\begin{proof}
The proof of
Theorem~\ref{thm:utp-common-arrow}
first constructs a canonical
matrix \(M\) from a genuine
full-Galois lift, by tame
inertia and the explicit
parabolic operations.
Before making the final
Tate-coefficient comparison,
this \(M\) already sends
\(1\) and every \(x_j/p^{k_j}\)
to \(I\setminus\beta I\).
If its lift is \(\alpha_0\),
Proposition~\ref{prop:uim-transfer}
shows that \(\alpha_0^{-1}\)
has canonical core action \(M\).
This proves
\eqref{eq:uim-common-minimum}
without a coefficient rescaling.

Every canonical core operator
preserves \(I\), hence each
\(p^kI\). Since \(1\in I\)
and \(x_j\in p^{k_j}I\),
each field component of its
tensor image has valuation
at least \(k_j-m_j\kappa\).
The component embeddings
preserve the native valuation,
so every such image belongs
to \(P_j\). The common
operator just constructed
gives this exact valuation
in every component. Its
tensor is therefore a
generator of \(P_j\) times
a unit of \(B_{m_j}\).
This gives the reverse
inclusion before closure;
the fractional ideal is closed.
Finally
Theorem~\ref{thm:uim-base-membership}
places all these orbit points
in the inverse image of the
raw union, proving
\eqref{eq:uim-lower-inclusion}.
\end{proof}

The quantifier order remains
\(\forall(\zeta_j)\,\exists F\,\forall j\).
Nothing requires one operator
to work for all torsion tuples.
Equation~\eqref{eq:uim-lower-inclusion}
does not identify the whole
possible-image hull with \(P_j\).
Other allowed representatives
and vertical branches may
enlarge it.

The difference between the
pure log-field point and a
one-plus-unit Kummer input
can also be detected exactly.

\begin{lemma}[Trace distinguishes two points with the same ideal behavior]
\label{lem:uim-trace-separation}
For \(s=j^2\), \(\zeta\in\mu_{2\ell}\), put
\[
 a=r_0^s,\quad
 k=\lfloor2s/\ell+\kappa\rfloor,\quad
 t=p^{-1}\log(1+p\zeta a),\quad
 \varepsilon=\zeta^\ell\in\{1,-1\}.
\]
Then
\begin{equation}\label{eq:uim-traces}
 \begin{aligned}
 \operatorname{Tr}_{E/\Q_p}(\zeta a/p^k)&=0,\\
 \operatorname{Tr}_{E/\Q_p}(t/p^k)
     &=\frac{30}{p^{k+1}}
             \log(1+\varepsilon p^\ell b_0^s)\ne0.
 \end{aligned}
\end{equation}
The valuation of the second line is
\(\ell-1+2s-k\ge\ell\).
\end{lemma}
\begin{proof}
The tame character of \(r_0^s\)
has order \(\ell\), because
\(\ell\nmid j^2\).
Its \(\ell\) conjugates sum
to zero, each occurring
fifteen times in \(E/K_0\);
\(\zeta\) is fixed by this
inertia. This proves the
first trace equality.
Multiplication of the same
conjugates gives
\[
 N_{E/K_0}(1+p\zeta r_0^s)
          =(1+\zeta^\ell p^\ell b_0^s)^{15}.
\]
The sign is positive because
\(\ell\) is odd. The right
side lies in \(\Q_p\),
so its norm from \(K_0\)
squares it. Taking logarithms
on principal units and
dividing by \(p^{k+1}\)
proves the second line of
\eqref{eq:uim-traces}.
Since \(p\nmid30\), its
first nonzero logarithmic
term gives valuation
\(\ell+2s-k-1\).
Finally \(k=\lfloor2s/\ell\rfloor+1\le s\)
for \(s\ge1\) and \(\ell\ge7\).
The lower bound follows.
\end{proof}

Allowed integral-coefficient
Galois/Kummer maps preserve
the trace-zero kernel by
Proposition~\ref{prop:trace-transport}.
Canonical core maps differ
from such maps by a unit
normalization, so they
preserve it as well.
The two points in
\eqref{eq:uim-traces} cannot
literally lie in the same
such orbit. This is consistent
with the per-operator
principal-ideal equality of
Proposition~\ref{prop:utp-logarithmic-bridge}.
The membership theorem used
the pure points from the
actual log-field pullback,
not a substitution of the
one-plus-log points for them.

The remaining boundary is global.
The initial curve, prime
\(\ell\), places and \(q\)
were fixed at the start.
Before Ind1, the copy-identity
preserves the root equation
and its marked value-group
generator. After Ind1, however,
one cannot conclude that
\(F\) fixes the native scalar
\(q\), the native ring unit
or the old global field
embedding. The transported
multiplicative data retain
their abstract value-group
generator; its valuation in
the old native chart is a
different datum.

The q-pilot in
\cite[Corollary~3.12, p.~174]{IUT3}
is explicitly left outside
Ind1, Ind2 and Ind3.
Keeping that original object
in the notation is not a
proof of the full horizontal
compatibility assertion.
Theorem~\ref{thm:uim-base-membership}
concerns a fixed base branch
in one column, not a transport
of that possible-image set
through a horizontal
theta-link.

One still has to match the
complete globally synchronized
pilot family, including the
other places and all required
vertical branches, with the
precise holomorphic reference,
global weights and region
used in the subsequent upper
bound. The local inclusion
does not identify the
one-prime bundle of
Theorem~\ref{thm:ub-realization}
with that full pilot.
Nor does it establish the
restriction to prime-strip-interpretable
regions made later in the
proof on p.~175 of
\cite{IUT3}.
The reconstruction comparisons,
actual full-Galois action
and source membership in
this section are mathematical
arguments using the specified
external reconstruction
results; they are not new
Lean formalizations.
No global \(abc\) conclusion
is inferred.
